\chapter{Exempelsaker}
Lekplats för att visa hur latex funkar.
Här är ett citat \cite{Torre2008}
haj haj inte blblblblblbl
khjjhg
nu ändrade kalle lite här

För att sätta en label på en ekvation behövs
\begin{verbatim}
\begin{equation}
\end{verbatim}
\begin{equation} \label{eq:viktig}
X \cdot 5 = \unit[4]{kg}
\end{equation}

Annars räcker det med två dollartecken

$$\theta=\alpha_2-\frac{3\pi}{2}+\alpha_1$$

Här är flera \cite{Torre2008}.
hej heje hejjejejej sami här allt bra eller aa det är bra här länge sen vi sågs nu här har du lite matte 
jag gillar $\unitfrac[3.40]{m}{s^2}$ eller en vanlig enhet $\unit[3.2]{N}$.
\begin{itemize}
	\item en lista
	\item trolol
\end{itemize}


\begin{enumerate}
	\item ett
	\item nr. två
\end{enumerate}


\section{tja}
Här refererar jag en ekvation \eqref{eq:viktig}. Och här är en ekvation $ \frac{a}{b} $ mitt inne i texten.


\begin{figure}[H]
\includegraphics[width=0.5\textwidth]{bitcoin}
\caption{Hej en viktig bild}
\label{pic:omgbild}
\end{figure}


